% !TEX root = ../main.tex
\clearpage
\phantomsection
\section*{6 \textbf{연구의 한계}}
\label{sec:limitations}
\addcontentsline{toc}{section}{6 연구의 한계}

본 연구는 EndorseRank를 Adaptive Weighted PageRank에 대해 엄밀히 평가하도록 설계되었으나, 다음과 같은 고유한 한계와 타당성 위협을 인정한다.


\proposalSubheading{1. 데이터 범위 및 품질}

\begin{itemize}
\item[\textbullet] \textbf{이벤트 완전성:} Arbitrum One 아카이브 또는 BigQuery의 공개 이벤트 로그(Approve, Permit, Transfer, GMX \textbf{PositionDecrease})에 의존한다. 체인 재구성 또는 노드 동기화 문제로 인한 누락·오인덱싱 이벤트가 그래프 구성 또는 GMX 결과 집계를 편향시킬 수 있다.
\item[\textbullet] \textbf{대리변수 한계:} 거래 성공 대리변수(청산 성공 횟수, 실현 이익 대리변수)는 도메인 직관 측정치이며 외부 ground-truth 신용 레이블이 아니다. 높은 상관만으로는 선행 대출 중심 점수 연구(예: Ghosh et al., 2024)의 무관한 담보 부족 대출 결과에 대한 예측 타당성을 확립하지 못한다.
\end{itemize}
\proposalSubheading{2. 시간적 범위}

\begin{itemize}
\item[\textbullet] \textbf{고정 기간:} 단일 6개월 기간으로 분석을 제한하면 장기 평판 역학 또는 시장 체제 변화를 포착하지 못할 수 있다. 결과는 특히 고변동성 또는 GMX 프로토콜 업그레이드 기간 등 다른 기간에서 달라질 수 있다.
\end{itemize}
\proposalSubheading{3. 그래프 모델링 가정}

\begin{itemize}
\item[\textbullet] \textbf{승인 의미론:} 0이 아닌 approve/permit이 실질적 신뢰 보증을 나타낸다고 가정한다. 실무에서는 사용자가 기본적으로 ``전액 지출'' 승인을 발행하거나 수탁 서비스를 통해 상호작용할 수 있어 신뢰 신호를 희석시킨다.
\item[\textbullet] \textbf{엣지 가중치 집계:} 모든 송금 합산 또는 최신 승인만 사용하면 자금 출처를 무시하고 다중 토큰 또는 다중 체인 상호작용을 가릴 수 있다.
\item[\textbullet] \textbf{GMX 거래 의미론:} 수익성 \textbf{PositionDecrease} 청산은 무기한 스왑 포지션의 실현 PnL을 반영하며 대출 상환이 아니다. 신용 유사성은 해석적이며, wash trading, copy trading, 자동화 전략이 진정한 신뢰성 없이 대리변수를 부풀릴 수 있다.
\end{itemize}
\proposalSubheading{4. 알고리즘적 단순화}

\begin{itemize}
\item[\textbullet] \textbf{고정 감쇠 인자:} 단일 감쇠 인자(d = 0.85)를 일괄 적용한다; 실제로는 토큰, 지갑 유형, 네트워크 세그먼트별로 최적 인자가 다를 수 있다.
\item[\textbullet] \textbf{모델링되지 않은 속성:} EndorseRank는 오프체인 요인(예: KYC 상태, IP 지리) 및 비 ERC-20 관계(예: NFT 승인, 크로스체인 브리지)를 무시하며, 이들도 신뢰에 영향을 줄 수 있다.
\item[\textbullet] \textbf{최소 청산 필터:} $\text{total\_closes} \geq 3$ 요건은 초보 또는 비활성 트레이더를 제외하고 매칭 표본에 선택 편향을 도입할 수 있다.
\end{itemize}
\proposalSubheading{5. 일반화 가능성}

\begin{itemize}
\item[\textbullet] \textbf{프로토콜 특수성:} Arbitrum One의 GMX V2 무기한 스왑 결과에 집중한다. EndorseRank 정렬은 다른 무기한 거래소(예: dYdX, Hyperliquid) 또는 비거래 사용 사례에서 달라질 수 있다.
\item[\textbullet] \textbf{자산 편향:} 고거래량 토큰(예: USDC, WETH)이 그래프 구조를 지배한다; 소규모 또는 신규 토큰은 PageRank에 덜 적합한 희소 보증 네트워크를 산출할 수 있다.
\end{itemize}
\proposalSubheading{6. 적대적 조작 및 Sybil 공격}

\begin{itemize}
\item[\textbullet] \textbf{그래프 게이밍:} 악의적 행위자가 더미 승인을 발행하거나 순환 승인 패턴을 조직하여 EndorseRank 점수를 부풀릴 수 있다. 유사하게 조율된 소규모 수익성 청산이 청산 성공 횟수를 부풀릴 수 있다. Sybil 전략 탐지 및 완화는 본 연구 범위를 넘어선 추가 안전장치를 요구한다.
\end{itemize}
\proposalSubheading{7. 계산 제약}

\begin{itemize}
\item[\textbullet] \textbf{확장성 한계:} EndorseRank가 표본에서 AWP보다 효율적이더라도, 전체 Arbitrum 송금 그래프(수백만 지갑)에서 실시간 계산은 여전히 분산 컴퓨팅 또는 근사 기법을 요구할 수 있다.
\end{itemize}
